\chapter{支持向量机与核方法}
\chaptergoal{通过几何图理解分类间隔和支持向量；区分线性与核方法；解释软间隔中的 $C$ 及 RBF 核中的 $\gamma$；使用标准化的 sklearn 流程进行有限参数比较，并了解 SVM 的计算与解释限制。}

支持向量机（support vector machine, SVM）展示了另一种分类思想：不只是找到一条能够区分类别的边界，还试图在合理的约束下形成较大的分类间隔。它的核心几何直觉相对容易理解，复杂优化推导可以放在选读部分。这一章安排较紧凑，以图形和实践为主。
\section{分类间隔与支持向量}

\textbf{分类间隔}可以先理解为分界线与附近训练样本之间留出的空间。考虑一个人工二维例子：一类点的横坐标为$-2,-1$，另一类为$1,2$，纵坐标对当前分界无影响。直线$x_1=0$把两类分开，最近点到它的距离为1，两侧边缘线$x_1=-1$和$x_1=1$之间宽度为2。如果把边界移到$x_1=0.8$，仍可能没有训练分类错误，但右侧最近点只离边界0.2；它对位置扰动更敏感。

\textbf{支持向量}是在SVM求解中对边界产生直接约束作用的训练记录。在这个理想可分例子里，边缘上的记录起关键作用；后面允许间隔违规时，支持向量还可能位于间隔内部，甚至被错分，所以不能把所有SVM的支持向量都定义成恰好压在两条边缘线上的点。

在二维线性可分的例子中，不同直线可能都能分开两类点。SVM 关注的是边界与最接近训练点之间的间隔，试图选择具有较大几何间隔的边界。在严格可分的最优间隔构造中，边缘上的关键样本构成支持向量。直观上，边界的位置由较少的关键训练点决定，而不是简单地对所有点求平均。

分类得分可以写成 $f(\mathbf x)=\mathbf w^\mathsf T\mathbf x+b$。类别输出由得分符号决定；$\|\mathbf w\|$ 参与定义几何间隔。这里的得分不是天然经过概率校准的类别概率。若在 sklearn 中启用概率输出，还会增加额外计算并需要理解其校准方式。

\fig{ch11_svm_margin.pdf}{二维生成数据上的线性间隔和支持向量示意，说明边界附近样本如何影响模型。}

\subsection{从得分到几何距离}
同一条分界可以有许多得分表达。例如 $f(x)=x$ 与 $f(x)=10x$ 的正负号相同，边界都在0，但得分大小相差10倍。若直接以“样本得分很大”作为间隔宽的依据，任意放大权重就能伪造大间隔。因此应除以法向量长度：点到线性边界的有符号几何距离为
\[
\frac{\mathbf w^\mathsf T\mathbf x+b}{\|\mathbf w\|_2}.
\]
对标签 $y_i\in\{-1,+1\}$，$y_i(\mathbf w^\mathsf T\mathbf x_i+b)/\|\mathbf w\|_2$ 为带类别方向的距离，正确侧为正，错误侧为负。将 $\mathbf w,b$ 同乘正数不改变它。

硬间隔形式通过规定所有训练记录满足 $y_i(\mathbf w^\mathsf T\mathbf x_i+b)\geq1$ 来固定得分尺度。在这个规范下，两条支持超平面 $f=1$ 与 $f=-1$ 间的宽度为 $2/\|\mathbf w\|_2$。因此最大化间隔等价于在这些约束下最小化 $\|\mathbf w\|_2^2/2$。这几步就建立了几何直觉与目标函数的关系，不必在正文展开完整对偶推导。

\worked{四个一维点的最大间隔}
取负类 $x=-2,-1$，正类 $x=1,2$。边界为 $wx+b=0$，在最近两个点上的约束是 $w-b\geq1$ 和 $w+b\geq1$，相加得到 $w\geq1$。选择 $w=1,b=0$ 达到最小范数，间隔宽度为2。两个较远点的约束也满足，但它们没有进一步缩小可行范围。

将 $x=1$ 改为0.4而保持正类，最近两类之间的空隙变小，最大间隔边界会移动。相反，把 $x=2$ 移到3，在这个简单硬间隔问题中不改变两个最近边界约束。支持向量的重要性由最优化条件体现，不是事先挑几个“看起来典型”的对象。

\subsection{硬间隔为什么需要谨慎使用}
若存在相同输入却有不同标签，不存在同时满足全部硬间隔约束的线性分类器。即使训练样本恰好可分，也可能是样本很少或某个测量量泄漏了目标；大间隔几何并不自动保证研究设计正确。实际科学分类常有类别重叠、误标和测量噪声，软间隔是更合理的起点。

“支持向量少”还不等于只有这些样本在科学上重要。它描述当前表示、核和参数下优化解的构成；换标准化、换核或改变数据，支持集合都可能改变。不能直接把支持向量当作独立的新天体族或物理临界状态。

\section{软间隔和参数\texorpdfstring{$C$}{C}}

\textbf{硬间隔}要求训练样本在规定的间隔外且分类正确；\textbf{软间隔}允许违反部分要求，并给这些违反增加代价。为写清这一点，暂把两类编码成$y\in\{-1,+1\}$，得分记为$f(\mathbf x)$。一种常见的间隔损失是 $\max(0,1-yf(\mathbf x))$，也叫铰链损失。

若真实类别为$+1$、得分为0.2，类别符号是对的，但间隔损失为0.8；若得分为$-0.2$，已经错分，损失为1.2；得分为1.5时损失为0。可见“违反间隔”不全等于“分类错误”。一个常见训练目标是 $\tfrac12\|\mathbf w\|^2+C\sum_i\max(0,1-y_if(\mathbf x_i))$：第一项控制边界的几何尺度，第二项累计间隔违规代价。这里只用它解释$C$的角色，不要求展开约束优化证明。\textsuperscript{[20]}

真实数据中的不同类别可能重叠，测量也可能有噪声，要求所有训练样本严格分开往往不现实。软间隔允许一些样本违反间隔条件，并用惩罚权衡间隔大小与训练错误。\texttt{SVC} 的参数 \texttt{C} 控制这一权衡：较大 \texttt{C} 通常更强调训练中的间隔违规惩罚，较小 \texttt{C} 更偏向允许一些违规以取得较宽松的边界。

需要避免机械记忆“C 大一定过拟合”。具体结果取决于特征、类别重叠、核函数与验证结构。训练误差和验证误差必须一起检查；若标签包含误标，过度强制拟合训练类别反而可能降低实际应用可靠性。

\subsection{铰链损失把两种违规区别开来}
记函数间隔 $u_i=y_if(\mathbf x_i)$。若 $u_i\geq1$，记录不仅分类正确，也在规定边缘外，铰链损失为0；若 $0<u_i<1$，类别正确但进入间隔，损失为 $1-u_i$；若 $u_i\leq0$，位于边界或错误侧，损失至少为1。三个区间解释了为什么SVM训练不只是统计0/1误分数。

引入非负松弛变量 $\xi_i$，条件可写成 $y_if(\mathbf x_i)\geq1-\xi_i$，目标为 $\|\mathbf w\|^2/2+C\sum_i\xi_i$。对给定得分，最小可行 $\xi_i=\max(0,1-y_if_i)$，于是得到前面的无约束铰链形式。两种写法描述同一权衡，不是两种不同分类算法。

\worked{比较两个都分类正确的候选}
仍用 $x=(-2,-1,1,2)$。候选A取 $w=1,b=0$，四条记录铰链损失全为0，目标为0.5。候选B取 $w=0.5,b=0$，两条内侧记录各损失0.5，外侧为0，所以目标为 $0.125+C$。

当 $C=0.1$，B的目标0.225小于A的0.5：B允许两条记录进入间隔，换取更小范数。若 $C=1$，B为1.125，A更好。这里仅比较两个候选，不声称全体可能参数的最优解只能是这两个，但已经说明 $C$ 控制的不是准确率公式中的直接权重，而是拟合违规与几何尺度的取舍。

本节目标采用损失\emph{求和}，若改成平均，惩罚系数需要随样本数对应调整。不能比较两个样本规模完全不同的实验时，仅因 $C$ 相同就认为正则化相对强度相同。不同库或模型使用相近参数名，也必须核对目标的归一化方式。

\subsection{类别代价、概率与科学使用}
少数类的漏检代价较高时，可以通过类别权重改变不同记录的间隔惩罚；这与训练后移动决策阈值有关，但并不等价，因为前者改变拟合出的边界，后者在固定得分上改变决策。应分别说明选择依据和验证方式。

SVM得分距离不能直接叫“属于正类的概率”。启用概率估计或另行校准会引入额外步骤。较大的正得分通常表示在模型判别方向上更偏向正类，但不同核、不同 $C$ 的数值尺度未必可直接比较。对候选筛选，排序指标与概率解释应当分开核验。

\section{核方法与非线性边界}

先看一维数据$x=-2,-1,1,2$，外侧两个点属于A，内侧两个点属于B。原坐标的一次阈值不能把两个外侧点同时放入同一个连续区间。然而构造新特征$z=x^2$后，A对应$z=4$，B对应$z=1$，用$z=2.5$作阈值就可以分开。这说明改变\textbf{特征表示}能够改变简单边界所能表达的关系。

\textbf{核函数}是一种满足相应数学条件的相似性函数，可等价地计算某个特征空间中的内积。例如显式变换$\phi(x)=x^2$对应$K(x,z)=x^2z^2$。一般核方法不必先把所有变换后的坐标列出来，而可以直接通过这些内积完成相关计算。不是任意称作“相似度”的函数都自动适合作为SVM的核。

当两个类别呈弯曲结构时，原始二维空间中的直线可能无法很好地分开样本。核方法通过样本间相似性的计算，等效利用更丰富的特征关系，而不必显式构造高维映射的全部坐标。常见 RBF 核为
\begin{equation}
K(\mathbf x,\mathbf z)=\exp\bigl(-\gamma\|\mathbf x-\mathbf z\|^2\bigr).
\end{equation}
核函数度量的是模型定义下的相似性，不一定等于物理空间中的几何距离。对于不同单位的输入，标准化尤其重要；否则某个尺度很大的变量可能主导核值。

图中的双月数据是程序构造的两条弯曲半环及其周围扰动，类别由生成过程指定，不是实测天体。scikit-learn的\texttt{make\_moons}函数用于生成这种机制演示数据。

\fig{ch11_svm_kernels.pdf}{人工生成的双月形数据上，线性核与 RBF 核的分类边界。用于说明表示能力差异，不对应真实物理类别。}

\subsection{一次可以手算的核技巧}
考虑二维向量 $\mathbf x=(x_1,x_2)$，显式构造三个二次特征
\[
\phi(\mathbf x)=(x_1^2,\sqrt2x_1x_2,x_2^2).
\]
直接计算内积有
\[
\phi(\mathbf x)^\mathsf T\phi(\mathbf z)
=x_1^2z_1^2+2x_1x_2z_1z_2+x_2^2z_2^2
=(\mathbf x^\mathsf T\mathbf z)^2.
\]
右侧只使用原始二维向量，却得到变换后三维内积。取 $\mathbf x=(1,2),\mathbf z=(2,1)$，右侧为 $4^2=16$；左侧为 $1\times4+(2\sqrt2)^2+4\times1=16$。配套程序核对两种计算的一致性。

核技巧的意义不是消除了非线性，而是避免显式存储某些高维特征。合法核在有限样本上形成的Gram矩阵具有相应的半正定性质。对刚才的映射，任意系数向量 $\mathbf a$ 有
\[
\sum_{ij}a_i a_jK(\mathbf x_i,\mathbf x_j)
=\left\|\sum_i a_i\phi(\mathbf x_i)\right\|^2\geq0,
\]
这说明为什么内积自然满足要求。正文不要求证明任意核的全部性质，但应知道随意定义一个“相似分数”不一定保持优化问题所需结构。

\subsection{核SVM怎样预测一个新对象}
核SVM的得分可表示为
\[
f(\mathbf x)=\sum_{i\in\mathrm{SV}}\alpha_i y_iK(\mathbf x_i,\mathbf x)+b.
\]
其中 $\alpha_i$ 是训练得到的系数，只有非零系数对应的记录进入支持集合SV。式子把新样本与支持记录逐一比较相似性，再按带符号权重组合；它不同于K近邻只取最近若干条记录的简单投票。\texttt{support\_vectors\_}、相关系数和类别顺序应按具体实现解释。

训练与预测成本也因此与样本规模及支持向量数量相关。核矩阵可能需要处理大量成对关系，预测一个样本也可能需要与许多支持向量计算核值。对很大的目录，可考虑线性模型、近似核特征或其他方法，而不是不加评估地把小数据设置放大。

\section{RBF 核与参数\texorpdfstring{$\gamma$}{gamma}}

RBF是radial basis function的缩写，本章使用其常见的高斯径向基核；它只通过两点之间的距离决定核值。取$\gamma=1$，距离为0时核值为1，距离为1时为$e^{-1}\approx0.3679$，距离为2时为$e^{-4}\approx0.0183$。把$\gamma$改为4，距离为1时就只剩$e^{-4}$，所以一个点对较远位置的相似性影响更快变小。这个数字例子比单独记忆“gamma越大越复杂”更能说明参数在做什么。

$\gamma$ 决定 RBF 相似性随距离衰减的速度。较大的 $\gamma$ 使局部相似性范围更窄，边界可能更加曲折；较小的 $\gamma$ 使相似性变化更平缓。它与 $C$ 共同决定模型的有效复杂度。在训练样本数量较大时，核 SVM 的训练与预测成本可能明显增长，因此不一定适合每一个百万量级目录任务。

\subsection{把核宽度写成可理解的距离}
若以 $\ell$ 表示特征空间中的长度尺度，可把RBF写成
\[
K(\mathbf x,\mathbf z)=\exp\left[-\frac{\|\mathbf x-\mathbf z\|^2}{2\ell^2}\right],
\qquad\gamma=\frac1{2\ell^2}.
\]
当距离为 $\ell$，核值为 $e^{-1/2}\approx0.6065$。$\ell$较大时更远记录仍较相似；$\gamma$较大则相反。这里的距离属于选定特征表示，不是自动对应星际距离或实验装置的空间距离。

若所有输入数值都乘 $a$，距离平方乘 $a^2$。要保持相同核值，需将 $\gamma$改为 $\gamma/a^2$。若仅改变某一列单位，就不能用单个全局 $\gamma$完全抵消，必须重新考虑各列尺度。这一计算解释了标准化和核参数之间的联系。

\subsection{两个极端帮助判断模型状态}
在固定、互不重合的样本上，$\gamma$很小时，许多核值接近1，各记录显得相似，模型可能难以区分局部结构；$\gamma$很大时，不同样本间核值趋近0，每个样本只在很局部起作用。后一种设置可以拟合训练细节，但对远离支持记录的新点，许多核项都很小，预测可能主要由截距等项决定。

这不是说“大 $\gamma$ 的训练准确率必定高”或“小 $\gamma$ 必定欠拟合”。$C$、样本分布和求解状态同时影响结果。诊断时比较训练和验证表现、支持向量比例，并观察输入是否落在开发样本覆盖范围内，比背诵单一参数口号更可靠。

\section{葡萄酒类别识别及参数实验}
\begin{lstlisting}[caption={Wine 中同时搜索核 SVM 的两个关键参数}]
from sklearn.datasets import load_wine
from sklearn.svm import SVC
from sklearn.preprocessing import StandardScaler
from sklearn.pipeline import Pipeline
from sklearn.model_selection import StratifiedKFold, GridSearchCV

X, y = load_wine(return_X_y=True)
from sklearn.model_selection import train_test_split
X, X_hold, y, y_hold = train_test_split(
    X, y, test_size=0.2, stratify=y, random_state=42
)
# All comparisons below use development X, y only.
pipe = Pipeline([("scale", StandardScaler()),
                 ("svc", SVC(kernel="rbf"))])
cv = StratifiedKFold(5, shuffle=True, random_state=7)
search = GridSearchCV(pipe,
                      {"svc__C": [0.1, 1, 10],
                       "svc__gamma": [0.01, 0.1, 1]},
                      cv=cv, scoring="accuracy")
search.fit(X, y)
print(search.best_params_, search.best_score_)
\end{lstlisting}

这段代码先留出保留数据，再只在开发部分搜索参数。搜索得到的最佳验证分数用于选择方案，不等于最终保留数据上的评价；本章不查看保留部分。

\subsection{读一个二维参数表，而不是只抄最好组合}
对142条葡萄酒开发记录，在相同五折内拟合标准化与RBF-SVM，得到以下平均准确率。每折标准化均只在训练部分拟合。
\begin{center}\small
\begin{tabular}{rrrr}\toprule
$C$&$\gamma=0.01$&$\gamma=0.1$&$\gamma=1$\\\midrule
0.1&0.4933&0.9653&0.4012\\
1&0.9791&0.9931&0.5554\\
10&0.9788&0.9860&0.5980\\\bottomrule
\end{tabular}\end{center}
\fig{v3_ch11_grid.pdf}{同一开发五折中的SVM参数对照。网格最高点是选择结果，不能直接当成未经选择的最终测试表现。}

本次 $C=1,\gamma=0.1$ 最高，均值约0.9931，折间标准差约0.0138。$\gamma=1$时三个 $C$ 的表现均明显较差，说明“越局部越灵活”不等于“在这些新样本上更有效”。$\gamma=0.01$从 $C=0.1$ 改为1有很大提升，再增到10几乎没有改变；这显示两个参数必须联合理解。

现在可以提出后续问题：低分组合是训练也差，还是只在验证差？哪些类别被混淆？标准化后的哪些特征控制了距离？这比围绕0.9931继续增加几十个附近候选值，更有利于理解问题。若继续细化网格，应承认这仍是开发过程，并保留原始搜索记录。

\subsection{与其他算法比较时怎样控制范围}
使用本书第6、7、9章的开发分割，分别为逻辑回归、近邻、森林和SVM采用合理预处理，保持特征与评价折相同。SVM相对某个模型的改善可能只涉及一两条记录，应结合逐折差异、错误对象和计算成本解释。某种方法分数稍低但训练快、模型易检查，也可能满足教学或科学应用需求。

\section{专业分类与 SVR 选读}
在 Sonar 或 SDSS 分类中，核方法应与逻辑回归、KNN 和树模型在一致条件下比较。Sonar 样本较少，折间波动可能较大，参数搜索过多容易让结果不稳定。专业应用还要考虑输入数据随仪器、噪声与观测条件改变时，核距离是否仍然有意义。

支持向量回归（SVR）使用类似的核方法处理连续目标，可以作为提高练习。

它常使用一个宽度参数$\epsilon$定义允许误差带，单个目标的损失为 $\max(0,|y-\hat y|-\epsilon)$。例如$\epsilon=0.5$时，误差0.3没有这部分损失，误差1.2的损失为0.7。这并不是宣布小误差在科学上不重要，而是模型所选择的训练惩罚形式；最终仍应以适当的MAE、RMSE及专业精度要求评价。
课堂必做仍以分类为主；进一步学习可研究优化问题的等价表达和最优条件；这些证明不作为本章必修内容。

\subsection{不敏感误差带与回归目标}
SVR用 $\epsilon$ 定义得分与目标之间不计这部分损失的误差带，可写为
\[
L_\epsilon(y,f)=\max(0,|y-f|-\epsilon).
\]
取三条误差 $(-0.2,0.7,-1.4)$，$\epsilon=0.5$时损失分别为0、0.2、0.9。它保留误差超过容忍带的线性惩罚，同时用权重范数控制函数复杂度。这与平方误差强调大残差的方式不同，不应把SVR直接解释为“用核做普通最小二乘”。

$\epsilon$具有目标单位。将目标从伏特改为毫伏而不改变 $\epsilon$，会改变训练所容忍的物理误差。若使用目标标准化，应在预测后逆变换，再用原单位评价，并把 $\epsilon$的含义换算清楚。$C$同样与目标和损失尺度有关。

\subsection{专业迁移的关键判断}
声呐数据用60个频带能量描述一条回波，邻近频带可能相关。标准化每列可以防止数值范围主导距离，但不会消除相关性；PCA或其他表示是否有帮助，必须纳入交叉验证再比较。天文颜色空间中的邻近也可能受到红移、消光和测量误差共同影响，核相似不等于同一物理身份。

当类别比例、仪器响应或观测信噪比改变时，应重新检查分类阈值和误差分布。SVM输出一个坚定的符号，并不能说明新对象处于训练覆盖范围。用模型作科学筛选，需要同时保留对象标识、输入质量、得分和后续核验信息。

\chapterpractice{先在人工双月数据上比较线性核与 RBF 核边界，再在 Wine 中对 \texttt{C} 和 \texttt{gamma} 做小型网格比较，记录结果和训练成本。若选 Sonar 作为专业迁移，报告不同划分下的分数波动。}
\chapterexercise
\begin{enumerate}
\item 什么是支持向量？它们为什么常出现在分类边界附近？
\item 为什么对 RBF 核 SVM 进行特征标准化通常十分重要？
\item 增大 $C$ 或 $\gamma$ 会产生什么直观影响？为何必须依赖验证判断？
\item 为什么核 SVM 在训练样本规模极大时可能面临明显的计算成本问题？
\end{enumerate}

\section*{深化练习与方法检验}
本组题强调解释、计算与验证；D编号在本章内独立计数。计算题应写出所用定义和中间步骤，实验题保留实际输出，包括不支持预期的结果。
\begin{enumerate}[label=D\arabic*.]
\item 一维二分类样本为 $x=-2,-1$ 属于负类，$x=1,2$ 属于正类。求一个规范化最大间隔分隔器，指出支持向量和两条间隔边界的距离。
\item 对样本的有符号得分 $yf(x)=1.5,0.4,-0.3$，分别计算 hinge loss。说明一个分类正确的样本是否一定没有损失。
\item 使用特征映射 $\phi(x_1,x_2)=(x_1^2,\sqrt2x_1x_2,x_2^2)$，证明内积等于 $(x^Tz)^2$。在 $x=(1,2),z=(3,4)$ 上数值核对。
\item RBF核的 $\gamma$ 变大时，一个固定非零距离对应的核值如何变化？为什么极大 $\gamma$ 可以造成训练表现很好而验证很差？
\end{enumerate}
